\documentclass[aspectratio=169,11pt]{beamer}
\usetheme{default}
\usefonttheme{professionalfonts}
\setbeamertemplate{navigation symbols}{}
\setbeamertemplate{footline}[frame number]
\setbeamertemplate{frametitle}{\vskip0.25em\insertframetitle\par\vskip0.15em\hrule}
\setlength{\tabcolsep}{3pt}
\renewcommand{\arraystretch}{1.12}
\usepackage[T1]{fontenc}
\usepackage[utf8]{inputenc}
\usepackage{lmodern}
\usepackage{booktabs,tabularx,ragged2e,makecell,adjustbox,capt-of,fontawesome5}
\usepackage{hyperref}
\hypersetup{colorlinks=true,allcolors=black}
\newcolumntype{Y}{>{\RaggedRight\arraybackslash}X}
\newcommand{\smalltable}{\scriptsize\setlength{\tabcolsep}{2.5pt}\renewcommand{\arraystretch}{1.08}}
\newcommand{\source}[1]{\vspace{2pt}{\tiny\textit{Source: #1}}}
\title{Dadra \& Nagar Haveli and Daman \& Diu}
\subtitle{Policy Review: Economic, Fiscal and Social Indicators}
\author{Ranveer Sharma}
\date{}
\begin{document}
\begin{frame}\titlepage\end{frame}

\begin{frame}{Overview}
This note consolidates the policy-relevant evidence base for the Union Territory of Dadra \& Nagar Haveli and Daman \& Diu (DNHDD), covering high-frequency indicators, industrial and fiscal statistics, labour-market and welfare outcomes, budgetary allocations, and identified evidence gaps.
\end{frame}

\section{High-Frequency Indicators}
\begin{frame}{1. High-Frequency Indicators (HFIs)}\smalltable
\begin{tabularx}{\textwidth}{p{0.19\textwidth}p{0.34\textwidth}Y}\toprule
Indicator & Latest verified observation & Interpretation boundary\\\midrule
Registered factory sector & 1,881 operating factories; Rs.33,072.36 crore GVA; 270,916 persons engaged, 2023--24 & ASI survey aggregate; nominal values; not whole-economy GVA\\
Unincorporated non-agricultural sector & 38,245 establishments; 86,610 workers; Rs.1,60,553 GVA per worker, 2023--24 & ASUSE small-UT estimate; inspect precision before subgroups\\
Merchandise exports & US\$4,645.68 million, FY 2024--25 & Gross exports, not local value added\\
Post-settlement SGST & Rs.1,316 crore, FY 2025--26 & Tax-accounting flow, not output\\
Consumer inflation & 6.22\%, August 2026 & Provisional DNHDD combined CPI; thin price sample\\
Electricity supplied & 10,852 MU, FY 2024--25; reported deficit nil & Supply balance does not measure interruptions or voltage quality\\
Banking & Rs.15,655 crore deposits and Rs.6,900 crore credit, March 2024 & RBI row label/geography after 2021 remains ambiguous\\
\bottomrule\end{tabularx}
\captionof{table}{Table 1: High-frequency and sectoral macro dashboard.}
\end{frame}

\section{Annual Survey of Industries}
\begin{frame}{2. Annual Survey of Industries (ASI)}\smalltable
\begin{tabularx}{\textwidth}{p{0.36\textwidth}rrr}\toprule
Measure & 2022--23 & 2023--24 & Nominal change\\\midrule
Operating factories & 1,804 & 1,881 & +4.27\%\\
Persons engaged & 271,246 & 270,916 & $-0.12\%$\\
Total emoluments & Rs.8,372.77 crore & Rs.9,149.69 crore & +9.28\%\\
Total output & Rs.2,10,900.04 crore & Rs.2,26,162.46 crore & +7.24\%\\
Gross value added & Rs.31,127.19 crore & Rs.33,072.36 crore & +6.25\%\\
Net value added & Rs.27,576.90 crore & Rs.28,995.93 crore & +5.15\%\\
Share of all-India ASI GVA & 1.42\% & 1.35\% & $-0.07$ pp\\
\bottomrule\end{tabularx}
\captionof{table}{Table 2: Annual Survey of Industries: registered factory sector.}
\end{frame}

\section{Parliamentary Evidence}
\begin{frame}{3. Selected Parliamentary Evidence (1/2)}\tiny\setlength{\tabcolsep}{2pt}\renewcommand{\arraystretch}{1.1}
\begin{tabularx}{\textwidth}{p{0.22\textwidth}p{0.42\textwidth}Y}\toprule
Question/date & Official observation & Analytical use\\\midrule
LS USQ 5483, 3 April 2025 & 76 DAJGUA villages; 6,944 ST households in kuccha houses; 8 villages without all-weather roads; 7 without JJM piped water & Targets spatial and tribal service gaps; reconcile vintages\\
LS USQ 1010, 2 December 2024 & 11,039 PMKVY, 12,202 JSS and 3,643 NAPS participants over differing periods & Training volume is not placement or job retention\\
RS USQ 445, 24 July 2023 & 13,877 Udyam registrations; 24 shutdown flags; 2,391 women-owned & Administrative stock, not active-enterprise census\\
RS USQ 2482, 13 March 2026 & Four PMFME-formalised enterprises by end-2025 & Implementation reach requires diagnosis\\
RS USQ 2618, 24 March 2025 & Rs.30 crore AMRUT 2.0 commitment; Rs.6 crore released/sanctioned & Separate commitment, release, expenditure and completion\\
LS USQ 1352, 3 December 2024 & Rs.135.17 crore PMMSY projects approved; 10 direct beneficiaries reported & Infrastructure package must not be divided by beneficiary count\\
RS USQ 639, 25 July 2023 & Rs.259 crore current DISCOM dues & Power adequacy and sector liquidity are distinct\\
\bottomrule\end{tabularx}
\captionof{table}{Table 3: Selected parliamentary implementation evidence.}
\end{frame}

\begin{frame}{4. Corporate Social Responsibility (CSR)}\smalltable
\begin{tabularx}{\textwidth}{p{0.25\textwidth}p{0.18\textwidth}Y}\toprule
Cut & Reported spend & Share / interpretation\\\midrule
Education & Rs.95.70 crore & 37.0\% of reported total; largest development-sector label\\
Health care & Rs.54.24 crore & 21.0\%; separate from poverty/hunger/malnutrition labels\\
Environmental sustainability & Rs.24.80 crore & 9.6\%; taxonomy is source-reported and not consolidated by fuzzy rules\\
Five largest normalised company contributors & Rs.79.87 crore & 30.9\%; Bhilosa Industries, JBF Industries, Time Technoplast, Savita Oil Technologies and Jai Corp\\
District not classified elsewhere & Rs.113.78 crore & 44.0\%; prevents reliable district incidence analysis\\
\bottomrule\end{tabularx}
\captionof{table}{Table 4: CSR spending profile in the supplied transaction extract.}
\end{frame}

\section{Labour, Budget and Finance}
\begin{frame}{5. PLFS 2023--24}\smalltable
\begin{tabularx}{\textwidth}{p{0.43\textwidth}rrrr}\toprule
Indicator & Male & Female & Persons & India\\\midrule
LFPR, age 15+ (\%) & 87.8 & 46.5 & 69.4 & 60.1\\
WPR, age 15+ (\%) & 85.3 & 46.0 & 67.8 & 58.2\\
Unemployment rate, age 15+ (\%) & 2.9 & 1.1 & 2.3 & 3.2\\
Youth unemployment, age 15--29 (\%) & -- & -- & 6.6 & 10.2\\
\bottomrule\end{tabularx}
\captionof{table}{Table 5: PLFS 2023--24 labour-market indicators.}
\end{frame}

\begin{frame}{6. Demand for Grants (Demand No. 54)}\smalltable
\begin{tabularx}{\textwidth}{p{0.34\textwidth}rrr}\toprule
Year & Revenue (Rs. cr, est.) & Capital (Rs. cr, est.) & Total (Rs. cr)\\\midrule
2024--25 Actual & 1,632 & 1,019 & 2,651\\
2025--26 BE & 1,706 & 1,076 & 2,782\\
2025--26 RE & 1,674 & 1,069 & 2,743\\
2026--27 BE & 1,731 & 1,104 & 2,835\\
\bottomrule\end{tabularx}
\captionof{table}{Table 6: Demand 54 expenditure remains capital-intensive. Totals as reported; revenue/capital split estimated from the published chart.}
\source{MHA, Detailed Demands for Grants 2026--27, Demand 54, \url{https://www.indiabudget.gov.in/doc/eb/sbe54.pdf}.}
\end{frame}

\begin{frame}{7. Fiscal and Financial Indicators}\tiny\setlength{\tabcolsep}{2pt}\renewcommand{\arraystretch}{1.1}
\begin{tabularx}{\textwidth}{p{0.24\textwidth}p{0.35\textwidth}Y}\toprule
Indicator & Observation & Use and caveat\\\midrule
Direct tax collections & Latest state-code series: Daman and Diu Rs.440.48 crore in FY2023--24 and Rs.998.03 crore in FY2024--25; DNH row blank after FY2021--22 & Coding break prevents a continuous merged-UT series; no state-level personal/corporate split\\
Scheduled bank deposits & Rs.9,979 crore March 2020; Rs.11,586 crore March 2021 & RBI table explicitly combines the two former UTs for these years\\
Scheduled bank credit & Rs.4,458 crore March 2020; Rs.4,089 crore March 2021 & Place-of-sanction concept; merged coverage explicit only for these years\\
CPI combined inflation & 2.58\% January 2026; 3.99\% May; 6.22\% August & Provisional monthly estimates; interpret acceleration cautiously\\
\bottomrule\end{tabularx}
\captionof{table}{Table 7: Selected fiscal and financial indicators.}
\end{frame}

\begin{frame}{8. VAHAN Registration Flows}\smalltable
\begin{tabularx}{\textwidth}{p{0.25\textwidth}rrr}\toprule
Year & All vehicles & Selected two-wheeler bundle & Motor cars\\\midrule
2019 & 11,237 & 5,071 & 3,823\\ 2020 & 11,767 & 7,729 & 2,900\\ 2021 & 14,884 & 8,733 & 3,533\\ 2022 & 18,287 & 9,736 & 4,245\\ 2023 & 21,858 & 10,548 & 5,001\\ 2024 & 23,997 & 12,008 & 5,820\\ 2025 & 27,148 & 14,258 & 6,947\\ 2026 to 21 September & 19,219 & 10,049 & 5,361\\\bottomrule
\end{tabularx}
\captionof{table}{Table 8: VAHAN registration flows.}
\end{frame}

\section{Living Standards and Sector Outcomes}
\begin{frame}{9. Vehicle Registration Tax}\smalltable
\begin{tabularx}{\textwidth}{p{0.22\textwidth}p{0.25\textwidth}p{0.28\textwidth}Y}\toprule
Vehicle and tax & DNHDD rate and tax & Maharashtra rate and tax & Difference\\\midrule
Non-diesel car & 4\%; Rs.40,000 & 11\%; Rs.1,10,000 & Rs.70,000 lower in DNHDD\\
Diesel car & 6\%; Rs.60,000 & 13\%; Rs.1,30,000 & Rs.70,000 lower in DNHDD\\
Motorcycle/tricycle at the same illustrative cost & 6\%; Rs.60,000 & 10--12\% by engine size; Rs.1,00,000--Rs.1,20,000 & Rs.40,000--Rs.60,000 lower\\
\bottomrule\end{tabularx}
\captionof{table}{Table 9: Illustrative one-time tax on a new Rs.10 lakh private vehicle.}
\end{frame}

\begin{frame}{10. MPCE 2023--24}\smalltable
\begin{tabularx}{\textwidth}{p{0.55\textwidth}r}\toprule
Category & Rs. per person per month\\\midrule
DNHDD rural & 4,450\\ India rural & 4,247\\ DNHDD urban & 6,876\\ India urban & 7,078\\\bottomrule
\end{tabularx}
\captionof{table}{Table 10: Monthly per-capita consumption expenditure, 2023--24.}
\source{MoSPI, HCES 2023--24 Fact Sheet, Statement 15.}
\end{frame}

\begin{frame}{11. Investment Friendliness Index (IFI) 2026}\smalltable
\begin{tabularx}{\textwidth}{p{0.55\textwidth}r}\toprule Pillar & Share of pillar maximum\\\midrule
Infrastructure & 36\%\\ Business climate & 14\%\\ Resources & 56\%\\ Government policy & 53\%\\ Regulatory ease & 30\%\\ Institutional environment & 50\%\\ Financial health & 38\%\\ Environment resilience & 32\%\\\bottomrule\end{tabularx}
\captionof{table}{Table 11: Investment Friendliness Index 2026: relative pillar performance. Values rounded.}
\source{NITI Aayog/CRISIL, Investment Friendliness Index 2026, p.91.}
\end{frame}

\begin{frame}{12. Multidimensional Poverty}\smalltable
\begin{tabularx}{\textwidth}{p{0.5\textwidth}rr}\toprule District & NFHS-4 (2015--16) & NFHS-5 (2019--21)\\\midrule Dadra and Nagar Haveli & 26.4\% & 10.8\%\\ Daman & 7.2\% & 6.6\%\\ Diu & 4.3\% & 1.6\%\\\bottomrule\end{tabularx}
\captionof{table}{Table 12: Multidimensional poverty declined, but DNH remains the outlier (headcount ratio, \%).}
\source{NITI Aayog, National MPI: A Progress Review 2023, p.282.}
\end{frame}

\begin{frame}{13. Health Indicators (Comparable)}\tiny\setlength{\tabcolsep}{2pt}\renewcommand{\arraystretch}{1.08}
\begin{tabularx}{\textwidth}{p{0.39\textwidth}rrY}\toprule Indicator & NFHS-6, 2023--24 & NFHS-5 & Change/read-through\\\midrule
Women who have ever used the internet & 78.4\% & 36.7\% & +41.7 pp\\ Women with 10 or more years of schooling & 44.8\% & 35.8\% & +9.0 pp\\ Women aged 20--24 married before 18 & 16.3\% & 26.4\% & $-10.1$ pp\\ First-trimester antenatal care & 89.4\% & 77.7\% & +11.7 pp\\ Institutional births & 96.3\% & 96.5\% & Broadly stable\\ Caesarean births & 28.6\% & 22.9\% & +5.7 pp; private facilities 43.6\%\\ Fully vaccinated, card or recall & 89.0\% & 94.9\% & Lower point estimate\\ Children under five stunted & 37.1\% & 39.4\% & Modest improvement; still high\\ Children under five wasted & 24.2\% & 21.6\% & Deterioration\\ Children under five underweight & 41.3\% & 38.7\% & Deterioration\\ Men with elevated BP or medication & 18.4\% & 15.4\% & Increase\\ Ever-married women reporting spousal violence & 19.9\% & 16.8\% & Increase; disclosure and sample error require care\\\bottomrule\end{tabularx}
\captionof{table}{Table 13: NFHS-6 priority indicators for DNHDD.}
\end{frame}

\begin{frame}{14. Education Indicators}\smalltable
\begin{tabularx}{\textwidth}{p{0.35\textwidth}r p{0.35\textwidth}r}\toprule Indicator & Value & Indicator & Value\\\midrule
Schools & 433 & Enrolment & 146,825\\ Teachers & 5,288 & PTR & 28\\ Zero-enrolment schools & 0 & Single-teacher schools & 0\\ Libraries & 432 & Playgrounds & 431\\ Electricity & 433 & Solar power & 95\\ Digital libraries & 6 & Kitchen gardens & 352\\\bottomrule\end{tabularx}
\captionof{table}{Table 14: UDISE+ school-system snapshot, 2024--25.}
\end{frame}

\section{Evidence Gaps and Reform}
\begin{frame}{15. Evidence Gaps (1/2)}\tiny\setlength{\tabcolsep}{2pt}\renewcommand{\arraystretch}{1.1}
\begin{tabularx}{\textwidth}{p{0.23\textwidth}p{0.32\textwidth}Y}\toprule Gap & Why it matters & Minimum viable product\\\midrule
GSDP/GVA & No defensible output, growth or sector-share story & Back-cast nominal/real series, methods, revisions and district indicators\\
Factory-sector depth & ASI now supplies UT aggregates but not a linked local productivity/resource panel & Annual sector panel with jobs, wages, capital, energy, water and pollution\\
IIP and high-frequency production & No DNHDD-specific published IIP with weights, revisions and methods & Monthly physical-production index\\
Credit and direct tax & Post-merger geography and RBI/CBDT concordance and documented merged-UT series personal/corporate or rural/urban splits are unclear & Concordance and documented merged-UT series\\
Migration, vacancies and wages & Cannot explain labour tightness or job quality & Quarterly employer survey plus worker panel\\
\bottomrule\end{tabularx}
\end{frame}

\begin{frame}{15. Evidence Gaps (2/2)}\tiny\setlength{\tabcolsep}{2pt}\renewcommand{\arraystretch}{1.1}
\begin{tabularx}{\textwidth}{p{0.23\textwidth}p{0.32\textwidth}Y}\toprule Gap & Why it matters & Minimum viable product\\\midrule
NFHS-6 microdata and learning & Fact sheet is current but subgroup precision and joint deprivation need microdata & Confidence intervals, district dashboard, MPI recomputation and learning outcomes\\
Water and pollution & Binding constraint lacks operational transparency & Metered water account and facility/cluster compliance dashboard\\
Tourism/fisheries value & Visits/regulation do not measure local value and ecological effects & Spend, stay, catch, price, cold-chain indicators\\
CSR/start-ups & CSR expenditure is available, but 44\% is district-unclassified and portal roles cannot be linked to companies & CIN/registration-keyed company, implementer, project and outcome datasets\\
\bottomrule\end{tabularx}
\captionof{table}{Table 15: Highest-priority evidence gaps.}
\end{frame}

\begin{frame}{16. Single Stop Data for Research}
A consolidated, longitudinal official-data workbook for Dadra \& Nagar Haveli and Daman \& Diu, compiled from Government of India and RBI sources, is maintained at:
\begin{center}
\vspace{0.35em}
\href{https://github.com/alwaysaccessible/dnhdd-longitudinal-official-data}{\fontsize{40}{44}\selectfont \faGithub}
\par\vspace{0.45em}
\href{https://github.com/alwaysaccessible/dnhdd-longitudinal-official-data}{\nolinkurl{https://github.com/alwaysaccessible/dnhdd-longitudinal-official-data}}
\end{center}
\end{frame}

\begin{frame}{17. Reform Proposals}\tiny\setlength{\tabcolsep}{2pt}\renewcommand{\arraystretch}{1.13}
\begin{tabularx}{\textwidth}{p{0.17\textwidth}p{0.46\textwidth}Y}\toprule Horizon & Actions & Success measure\\\midrule
0--12 months & Statistical task force; clearance process map; water accounts; project/asset register; EV policy; labour observatory design & Published baselines, owners, timelines and service levels\\
12--36 months & Cluster common services; worker housing pilots; apprenticeship outcomes; wastewater reuse; tourism observatory; district compacts & Use, conversion, reliability and beneficiary outcomes\\
3--5 years & Brownfield industrial renewal; supplier upgrading; integrated mobility; aquifer recovery; clean fleets; higher-value resilience tourism/fisheries & Productivity, value retained locally, resource intensity and resilience\\
\bottomrule\end{tabularx}
\captionof{table}{Table 16: Sequenced reform programme.}
\end{frame}

\begin{frame}{Closing note}
The review is designed as an evidence base for a DNHDD policy dashboard: preserve the distinctions between measured outcomes, administrative stocks, provisional estimates, and unresolved data gaps when using the indicators for policy decisions.
\end{frame}
\end{document}
